\documentclass[13pt, handout]{beamer}

\usetheme{Madrid}
\usecolortheme{beaver}
\setbeamertemplate{navigation symbols}{}
\setbeamertemplate{page number in head/foot}{}
\setbeamercolor*{item}{fg=red}

\usepackage[portuguese]{babel}
\usepackage{tikz}
\usepackage{ragged2e}

\setbeamertemplate{footline}
{
  \leavevmode%
  \hbox{%
    \begin{beamercolorbox}[wd=0.3\paperwidth,ht=2.5ex,dp=1ex,center]{author in head/foot}
      Grupo 01 (DIUM)
    \end{beamercolorbox}%
    \begin{beamercolorbox}[wd=0.4\paperwidth,ht=2.5ex,dp=1ex,center]{title in head/foot}
      Gestão de uma cadeia de restaurantes \textit{fast-food}
    \end{beamercolorbox}%
    \begin{beamercolorbox}[wd=0.3\paperwidth,ht=2.5ex,dp=1ex,center]{date in head/foot}
      \insertshortdate\hspace*{2ex}
    \end{beamercolorbox}%
  }%
}

\title{Gestão de uma cadeia de restaurantes \textit{fast-food}}
\subtitle{{}\color{black}Desenvolvimento de Sistemas de Software}
\author[Grupo 01]{
    \begin{tabular}{lc}
         André Santos & A106854 \\
         Daniel Parente & A107363 \\
         Pedro Ferreira & A107292 \\
         Tiago Martins & A106938
    \end{tabular}
}

\institute[DIUM]{
    Universidade do Minho -- Licenciatura em Engenharia Informática
}
\date[16/01/2026]{\scriptsize 16 de janeiro de 2026}

\begin{document}
\frame{\titlepage}

\begin{frame}{Índice}
    \tableofcontents
\end{frame}

\justifying

\section{Análise de Requisitos}

\begin{frame}{Análise de Requisitos -- O Problema}
    
Pretende-se desenvolver \textit{software} que:

\vspace{0.5cm}

\begin{itemize}
    \item Permita o registo e gestão de pedidos efetuados por clientes;
    \item Suporte a configuração livre do pedido por parte do cliente;
    \item Assegure a preparação completa dos pedidos antes da entrega;
    \item Garanta a gestores consultarem facilmente estatísticas sobre a cadeia de restaurantes.
\end{itemize}

\end{frame}

\begin{frame}{Análise de Requisitos -- Modelo de Domínio}

    \centering
    \includegraphics[width=0.85\textwidth]{dominio.png} \\
    {\scriptsize Modelo de Domínio construído.}

\end{frame}

\begin{frame}{Análise de Requisitos -- Identificação dos Casos de Uso}

        \centering
        
        \begin{minipage}{0.45\textwidth}
            \includegraphics[width=\textwidth]{geral.png}
        \end{minipage}
        \hspace{0.4cm}
        \begin{minipage}{0.45\textwidth}
            \includegraphics[width=\textwidth]{gestaodepedidos.png}
        \end{minipage} \\
        \vspace{0.25cm}
        {\scriptsize Diagramas de casos de uso construídos.}

        \begin{itemize}
            \item Identificação dos atores e das interações com o sistema
            \item Construção dos diagramas apropriados
        \end{itemize}
        
\end{frame}

\begin{frame}{Análise de Requisitos -- Especificação dos Casos de Uso}

    \centering
    \includegraphics[width=0.9\textwidth]{especificacao.png} \\
    {\scriptsize Exemplo de uma especificação de caso de uso -- UC Mandar Mensagem.}
    
\end{frame}

\section{Modelação Conceptual}

\begin{frame}{Modelação Conceptual -- Definição da API da Camada de Negócio}

    \centering
    \includegraphics[width=\textwidth]{api.png} \\
    {\scriptsize Definição da API da camada de negócio para o caso de uso "Preparar Etapa do Pedido".}

    \begin{itemize}
        \item Identificação das responsabilidades da camada de negócio para cada caso de uso;
        \item Definição da API.
    \end{itemize}
\end{frame}

\begin{frame}{Modelação Conceptual -- Diagramas de Sequência}

    \centering
    \includegraphics[width=0.7\textwidth]{sequencia.png} \\
    {\scriptsize Diagrama de sequência da operação de atualização do estado do pedido.}

    \vspace{0.5cm}
    \begin{itemize}
        \item Modelação comportamental de todas as operações da camada de negócio.
    \end{itemize}
    
\end{frame}

\begin{frame}{Modelação Conceptual -- Divisão por Subsistemas}

    \centering
    \includegraphics[width=0.7\textwidth]{componentessemdaos.png} \\
    {\scriptsize Primeiro diagrama de componentes construído, sem DAOs.}
    
\end{frame}

\begin{frame}{Modelação Conceptual -- Diagrama de Classes}

    \centering
    \includegraphics[width=1\textwidth]{subsystemcozinha.png} \\
    {\scriptsize Diagrama de classes do subsistema de cozinha.}
    
\end{frame}

\section{Implementação}

\begin{frame}{Implementação -- BD Relacional}

    \centering
    \includegraphics[width=1\textwidth]{camadadedados.png} \\
    {\scriptsize Hierarquia dos DAOs, utilizados para ORM.}
    
\end{frame}

\begin{frame}{Implementação -- Arquitetura Final}

    \centering

    \begin{minipage}{0.45\textwidth}
        \includegraphics[width=\textwidth]{componentes.png}
    \end{minipage}
    \hspace{0.5cm}
    \begin{minipage}{0.45\textwidth}
        \includegraphics[width=\textwidth]{packages.png}
    \end{minipage} \\
    \vspace{0.25cm}
    {\scriptsize Diagramas de componentes e \emph{packages} finais (com alguns dos DAOs representados).}
    
\end{frame}

\begin{frame}{Implementação -- Alterações aos Diagramas de Sequência}

    \centering
    \includegraphics[width=0.7\textwidth]{sequenciaDAOs.png} \\
    {\scriptsize Diagrama de sequência apresentado anteriormente, acomodado aos DAOs.}

    \vspace{0.5cm}
    
    \begin{itemize}
        \item Adição de atualizações à base de dados;
    \end{itemize}
    
\end{frame}

\begin{frame}{Implementação -- Camada de Apresentação}

    \begin{itemize}
        \item Composição de um manual de utilização para os utilizadores;
        \item Camada de apresentação na arquitetura MVC.
    \end{itemize}

    \vspace{0.5cm}

    \centering
    \includegraphics[width=0.7\textwidth]{view.png} \\
    {\scriptsize Menu que aparece após correr o programa.}
    
\end{frame}

\frame{\titlepage}

\end{document}
